\subsubsection{Đánh giá qua cross-validation}

\tablename~\ref{tab:classification-compare-models-cv-part1}
và~\ref{tab:classification-compare-models-cv-part2} thể hiện các kết quả trung
bình và độ lệch chuẩn các chỉ số Accuracy, Precision (macro), Recall (macro) và
F1 (macro) thông qua $5$-fold cross-validation.

\begin{table}[htbp]
  \centering
  \caption{Kết quả trung bình và độ lệch chuẩn các chỉ số của các mô hình
  thông qua cross-validation (phần 1)}
  \import{\tablespath}{classification-compare-models-cv-part1}
  \label{tab:classification-compare-models-cv-part1}
\end{table}

\begin{table}[htbp]
  \centering
  \caption{Kết quả trung bình và độ lệch chuẩn các chỉ số của các mô hình
  thông qua cross-validation (phần 2)}
  \import{\tablespath}{classification-compare-models-cv-part2}
  \label{tab:classification-compare-models-cv-part2}
\end{table}

Kết quả cross-validation trong
\tablename~\ref{tab:classification-compare-models-cv-part1}
và~\ref{tab:classification-compare-models-cv-part2} nhìn chung phù hợp và nhất
quán với kết quả trên tập kiểm tra, đồng thời cung cấp thêm thông tin về độ ổn
định của các mô hình thông qua độ lệch chuẩn.

Trước hết, Softmax tiếp tục duy trì Accuracy cao nhất (0.7122 ± 0.0009), rất
gần với kết quả trên tập test (0.7148). Độ lệch chuẩn nhỏ cho thấy mô hình rất
ổn định qua các lần chia dữ liệu. Tuy nhiên, Recall (0.4212) và F1 (0.4155) vẫn
thấp, củng cố nhận định trước đó rằng mô hình này thiên về tối ưu độ chính xác
tổng thể nhưng chưa xử lý tốt sự mất cân bằng giữa các lớp.

Weighted Softmax một lần nữa nổi bật với Recall cao nhất (0.6899 ± 0.0040), gần
như trùng khớp với kết quả test (0.697). Điều này cho thấy khả năng tổng quát
hóa tốt trong việc nhận diện các lớp, đặc biệt là lớp thiểu số. Tuy nhiên,
Accuracy (0.5903) và Precision (0.4556) vẫn thấp hơn Softmax, phản ánh sự đánh
đổi tương tự như đã quan sát trên tập kiểm tra. Độ lệch chuẩn nhỏ cho thấy mô
hình này cũng ổn định.

LDA tiếp tục thể hiện vai trò là mô hình cân bằng nhất, với các chỉ số Recall
(0.5791) và F1 (0.5108) khá tốt và ổn định. Kết quả này gần tương ứng với tập
test (Recall 0.5766, F1 0.5096), cho thấy LDA có khả năng tổng quát hóa đáng
tin cậy và ít bị dao động theo dữ liệu.

Trong khi đó, QDA vẫn có Accuracy rất thấp (0.2427) nhưng Recall tương đối cao
(0.5548), giống với kết quả trên tập test. Tuy nhiên, F1 thấp (0.2572) cho thấy
mô hình dự đoán thiếu chính xác và không cân bằng. Dù độ lệch chuẩn nhỏ (tức là
ổn định), nhưng đây là sự ổn định ở mức hiệu năng kém, nên không có giá trị
thực tiễn cao.

Tổng thể, kết quả cross-validation xác nhận lại các kết luận từ tập kiểm tra:
Softmax tốt nhất về Accuracy, Weighted Softmax tốt nhất về Recall, LDA cân bằng
nhất, và QDA kém hiệu quả. Đồng thời, độ lệch chuẩn nhỏ ở tất cả các mô hình
cho thấy các kết luận này đáng tin cậy và không phụ thuộc mạnh vào cách chia dữ
liệu.
